\documentclass[11pt]{article}
\usepackage{fontspec}
\usepackage[french,english]{babel}
\usepackage{amsmath}
\usepackage{amssymb}
\usepackage{amsthm}
\usepackage{mathrsfs}
\usepackage{tikz-cd}
\usepackage{enumitem}
\usetikzlibrary{babel}

\newcommand{\printedpage}[1]{\marginpar{\footnotesize p.~#1}}
\newcommand{\msp}[1]{\marginpar{\footnotesize [#1]}}
\newcommand{\Ind}{\operatorname{Ind}}
\newcommand{\Ob}{\operatorname{Ob}}
\newcommand{\Ord}{\operatorname{Ord}}
\newcommand{\card}{\operatorname{card}}
\newcommand{\cardred}{\mathop{\textup{card\,red}}\nolimits}
\newcommand{\cercle}[1]{\tikz[baseline=(c.base)]\node[draw,circle,inner sep=1pt](c){\textup{#1}};}

\begin{document}

\setcounter{footnote}{143}

\printedpage{129}
% head: A. Grothendieck : Dérivateurs, Ch. XVIII, Édition des Universités de Montpellier II et Paris VII
% foot: Version 17 décembre 2025

\msp{page 340}
with the exception of a$'$ $\Longrightarrow$ a, which is a little simpler in the present case. We are going to bring it out in the form of a somewhat more general lemma. (\textbf{N.B.} The proof of a$'$ $\Longrightarrow$ a in 4.13.15 had been omitted as ``sorital''\,!, cf. page 337\,\dots)

\textbf{Lemma 4.13.19.} \textit{Let $\mathcal{M}$ be an accessible category in $\mathcal{M}$, $\pi$ a useful cardinal for $\mathcal{M}$ and such that $\underline{S}$ is $\pi$-accessible, i.e. $\underline{S}$ is the essential image of $\Ind_\pi(\underline{S}(\pi))$ under the fully faithful canonical functor $\Ind_\pi(\underline{S}(\pi))\longrightarrow\mathcal{M}$, where $\underline{S}(\pi)=\underline{S}\cap\mathcal{M}_\pi$. We assume that there exists a limit ordinal $\alpha$, such that $\card\alpha\leqslant\pi$ and such that $\mathcal{M}_\pi$ and $\underline{S}(\pi)$ $(=\underline{S}_\pi)$ are stable under $\varinjlim$ of type $I_\alpha$ are stable under $\varinjlim$ of type $I_\alpha$, and that the inclusion functor $\underline{S}(\pi)\longrightarrow\mathcal{M}_\pi$ commutes with them. Let $\Phi$ be a set of categories such that $\forall\ I\in\Phi$, we have $\cardred I\leqslant\pi$, and suppose that $\mathcal{M}_\pi$ is stable under $\varinjlim$ of type $\Phi$, and that $\underline{S}(\pi)$ is stable in $\mathcal{M}_\pi$ under the said limits (or what amounts to the same thing, that it is stable in itself under these $\varinjlim$\,, and that the inclusion functor commutes with them). Then $\mathcal{M}$ and $\underline{S}$ are stable under $\varinjlim$ of type $\Phi$, and the inclusion functor $\underline{S}\longrightarrow\mathcal{M}$ commutes with them, i.e. $\mathcal{M}$ is stable under these limits, and $\underline{S}$ stable in $\mathcal{M}$ under these same limits.} % [sic] « soient stables par lim de type I_alpha » printed twice -> once, repetition

\msp{page 341}
\textsc{Proof}. We apply 4.11.1 (p. 267-268) separately to $\mathcal{M}\simeq\Ind_\pi\mathcal{M}_\pi$ and to $\underline{S}\simeq\Ind_\pi\underline{S}(\pi)$, to conclude that $\mathcal{M}$ and $\underline{S}$ are separately stable under $\varinjlim$ of type $\Phi$. It remains to see that the functor $\underline{S}\longrightarrow\mathcal{M}$ commutes with them. This is seen as in the proof of 4.11.1, more generally we get this:

\textbf{Lemma 4.13.20} (cor. to 4.11.1). \textit{Let $f_0 : C\longrightarrow C'$ be a functor between categories equivalent to small categories, $\pi$ an infinite cardinal, $f:\mathcal{M}\stackrel{\textup{def}}{=}\Ind_\pi C\longrightarrow\mathcal{M}'\stackrel{\textup{def}}{=}\Ind_\pi C'$ the canonical extension $\Ind_\pi(f_0)$ of $f_0$. Let $\Phi$ and $\alpha$ be as in 4.11.1, and suppose that $I_\alpha\in\Phi$ and that $f_0$ commutes with $\varinjlim$ of type $\Phi$. Then $f$ commutes with these same direct limits.} (\footnote{It is better to take functors $f_0 : C\longrightarrow\mathcal{N}$, and $\mathcal{N}$ [?] that $f_0$ commutes with $\varinjlim$ of type $\Phi$, so the same is true of its extension $\Ind_\pi(C,\mathcal{N})\longrightarrow\mathcal{N}$. (\textbf{N.B.} We assume that $\mathcal{N}$ satisfies $\mathrm{L}_\pi$.)}).

\textbf{Corollary 4.13.21.} \textit{Under the conditions of 4.13.20, suppose that $\Phi$ contains the categories of cardinality $\leqslant\pi$ (resp. the filtering categories of cardinality $\leqslant\pi$, resp. the filtering categories large relative to $\pi_0$ of cardinality $\leqslant\pi$, resp. the discrete categories of cardinality $\leqslant\pi$). Then $f$ commutes with small $\varinjlim$ (resp. with small filtering $\varinjlim$, resp. with small filtering $\varinjlim$ large relative to $\pi_0$, resp. with small sums). (Converse clear and ``left to the reader''.)}

Seen as in 4.4.14 and 4.4.14.1 (p. 97, 98).

\msp{page 342}
Here at last, it seems to me, is the definitive form (\,?) of 4.13.15, 4.13.18.

\textbf{Theorem 4.13.22} (\footnote{cf. variant 4.13.36, page 358.}). \textit{Let $\mathcal{M}$ be an accessible category, $\pi$ an infinite cardinal, useful for $\mathcal{M}$, and $\Phi$ a set of categories satisfying the two conditions:}
\begin{itemize}[leftmargin=2.5em]
\item[\textit{1°)}] $J\in\Phi$ \textit{implies} $\cardred J\leqslant\pi$.
\item[\textit{2°)}] \textit{There exists a} limit ordinal $\alpha_0$ \textit{with} $I_{\alpha_0}\in\Phi$.
\end{itemize}
\textit{We assume $\mathcal{M}$ stable under $\varinjlim$ of type $\Phi$. Let $\underline{S}$ be a} strictly full \textit{subcategory of $\mathcal{M}$, and $\underline{S}(\pi)=\underline{S}\cap\mathcal{M}_\pi$. Let us denote by $L(\varphi)$ the condition on $\underline{S}(\pi)$: $\underline{S}(\pi)$ is stable in} % [sic] L(\varphi) -> L(\Phi), the condition is later written L(\Phi)
\printedpage{130}
% head: A. Grothendieck : Dérivateurs, Ch. XVIII, Édition des Universités de Montpellier II et Paris VII
% foot: Version 17 décembre 2025
\textit{$\mathcal{M}$ (where what amounts to the same thing, in $\mathcal{M}_\pi$) under $\varinjlim$ of type $\Phi$. This being said, the following conditions on $\underline{S}$ are equivalent:} % [sic] où ce qui revient -> ou ce qui revient, conjunction "ou"
\begin{itemize}[leftmargin=2.5em]
\item[\cercle{a}] \textit{$\underline{S}$ is $\pi$-accessible in $\mathcal{M}$, and stable in $\mathcal{M}$ under $\varinjlim$ of type $\Phi$.}
\item[a$'$)] \textit{$\underline{S}$ is $\pi$-accessible in $\mathcal{M}$, and $\underline{S}(\pi)$ satisfies $L(\Phi)$.}
\end{itemize}
\msp{page 343}
\begin{itemize}[leftmargin=2.5em]
\item[b)] \textit{$\underline{S}(\pi)$ satisfies $L(\Phi)$, and for every $x\in\Ob\mathcal{M}$, we have this: $x\in\Ob\underline{S}$ if and only if the strictly full subcategory $\underline{S}(\pi)/x$ of $\mathcal{M}(\pi)/x$ is cofinal in $\mathcal{M}(\pi)/x$} \texttt{[read} $\mathcal{M}_\pi/x$\texttt{]}.
\item[\cercle{c}] \textit{$\underline{S}(\pi)$ satisfies $L(\Phi)$, and for every $x\in\Ob\mathcal{M}$ and every representation $x\simeq\varinjlim\limits_{i\in I}x_i$ which is $\Phi$-$\pi$-adapted (cf. definition below), we have this: $x\in\Ob\underline{S}$ if and only if $I_{\underline{S}}$ is cofinal in $I$, where $I_{\underline{S}}$ is the strictly full subcategory of $I$ formed by the $i\in\Ob I$ such that $x_i\in\Ob\underline{S}$ (or again, $x_i\in\Ob\underline{S}(\pi)$).} \textbf{N.B.} \textit{We say that a direct system $(x_i)_{i\in I}$ in $\mathcal{M}$ is $\Phi$-$\pi$-adapted if}
\begin{itemize}[leftmargin=2em]
\item[\textit{1)}] \textit{$I$ equivalent to a small category, and $I$ large relative to $\pi$,}
\item[\textit{2)}] \textit{the $\varinjlim$ of type $\Phi$ exist in $I$,}
\item[\textit{3)}] \textit{the functor $i\longmapsto x_i$ from $I$ to $\mathcal{M}$ commutes with $\varinjlim$ of type $\Phi$, and finally}
\item[\textit{4)}] \textit{the $x_i\in\Ob\mathcal{M}_\pi$.}
\end{itemize}
\item[c$'$)] \textit{(If $\pi^{\pi_0}=\pi$, where $\pi_0$ is a cardinal $\leqslant\pi$ useful for $\mathcal{M}$, and if $J\in\Phi$ implies that $J$ is large relative to $\pi_0$.) As in c), with $I$} ordered.
\end{itemize}
(\footnote{\textbf{N.B.} When the restrictive conditions opposite \texttt{[parenthesis in c$'$)]} are not satisfied, then a, a$'$, b, c are still equivalent, and imply c$'$, but are perhaps not equivalent to c$'$).},\footnote{\textbf{N.B.} One should restate the auxiliary conditions [?] (If \dots), cf. p. 359-360.}).

\textbf{Definition 4.13.23} (\footnote{\textbf{N.B.} It would be worthwhile to move this definition up (\textit{without} b), even before 4.13.15.}). Let $\mathcal{M}$ be an accessible category, $\pi$ a cardinal, $\Phi$ a set of categories, satisfying the two conditions:
\begin{itemize}[leftmargin=2.5em]
\item[a)] $J\in\Phi\Longrightarrow\cardred J\leqslant\pi$.
\item[b)] There exists a limit ordinal $\alpha_0$, such that $I_{\alpha_0}\in\Phi$ (\footnote{If it is to give a definition, one can dispense with condition b).}).
\end{itemize}
\msp{page 344}
Let $S$ be a subset of $\Ob\mathcal{M}$, $\underline{S}$ the full subcategory it defines. We say that $S$ is a \textit{$\Phi$-$\pi$-accessible subset of} $\Ob\mathcal{M}$, or that $\underline{S}$ is a \textit{$\Phi$-$\pi$-accessible subcategory in} $\mathcal{M}$, if the following conditions are satisfied.
\begin{itemize}[leftmargin=2.5em]
\item[1°)] $\pi$ is useful for $\mathcal{M}$, $\mathcal{M}$ is stable under $\varinjlim$ of type $\Phi$ (\footnote{Preliminary condition on $\mathcal{M}$.}).
\item[2°)] $\underline{S}$ is a $\pi$-accessible subcategory in $\mathcal{M}$.
\item[3°)] $\underline{S}$ is stable in $\mathcal{M}$ under $\varinjlim$ of type $\Phi$.
\end{itemize}
(\footnote{\textbf{N.B.} When $\Phi=\{I_{\alpha_0}\}$, we also say that $\underline{S}$ is $\alpha_0$-$\pi$-accessible in $\mathcal{M}$. If $\Phi$ consists of the categories of cardinality $\leqslant\pi$ and large relative to $\pi_0$, we say that $\underline{S}$ is $\pi_0$-$\pi$-accessible in $\mathcal{M}$.}).

\textbf{N.B.} When the preliminary condition 1°) \textit{on} $\mathcal{M}$ is satisfied, the validity of conditions 2°) 3°) is therefore equivalent to the condition that $\underline{S}$ be \textit{strictly} full in $\mathcal{M}$, and that in
\printedpage{131}
% head: A. Grothendieck : Dérivateurs, Ch. XVIII, Édition des Universités de Montpellier II et Paris VII
% foot: Version 17 décembre 2025
addition $\underline{S}$ satisfy condition a) of theorem 4.13.22. On the other hand, this condition is equivalent to each of conditions a$'$, b, c of this same theorem, and also to the last condition c$'$, provided that we have $\pi^{\pi_0}=\pi$ (for a cardinal $\pi_0<\pi$ useful for $\mathcal{M}$) and that $J\in\Phi$ implies that $J$ is large relative to $\pi_0$.

Conceptually, it is the definition given in 4.13.23, i.e. condition \cercle{a} of 4.13.22, which seems the clearest, or its immediate variant a$'$. But from the practical point of view, of working with accessible subsets, it is by property \cercle{c} which seems the most convenient. The main weight of theorem 4.13.22 is indeed
\msp{page 345}
the implication a $\Longrightarrow$ c (via a $\Longrightarrow$ a$'$ $\overset{!}{\Longrightarrow}$ c). So one must above all see this theorem as giving a remarkable criterion for membership in $\Ob\underline{S}$ of an object $x$ of $\mathcal{M}$ which is presented by a $\Phi$-$\pi$-adapted direct system:
\[ x\ \in\ \Ob\underline{S}\qquad\quad\Longleftrightarrow\qquad\quad I_{\underline{S}}\ \text{cofinal}\ , \tag*{$(*)$} \]
and in this equivalence, it is the implication $\Longrightarrow$ which is the non-sorital part, and remarkable.

\textbf{Proposition 4.13.24.} \textit{Let $\mathcal{M}$, $\pi$, $\Phi$ be as in 4.13.23} (\footnote{\textbf{N.B.} The condition $I_{\alpha_0}\in\Phi$ is not needed.})\textit{, we assume that condition 1° of 4.13.23 is satisfied ($\pi$ useful for $\mathcal{M}$, $\mathcal{M}$ stable under $\varinjlim$ of type $\Phi$). Let ${\Sigma'}^{\Phi}_{\pi}$ be the set of subcategories $\underline{S}$ of $\mathcal{M}$ which are $\Phi$-$\pi$-accessible in $\mathcal{M}$, and $\Sigma^{\Phi}_{\pi}$ the set of subcategories $\underline{S}(\pi)$ of $\mathcal{M}_\pi$ which are strictly full, and stable in $\mathcal{M}_\pi$ under $\varinjlim$ of type $\Phi$. (These subcategories are necessarily Karoubian, since stable under $\varinjlim$ of type $I_{\alpha_0}$, hence also under direct factors.) Then the bijections (4.13.4.1) (page 288) induce mutually inverse bijections}
\[
\begin{tikzcd}
\Sigma^{\Phi}_{\pi} \arrow[r, shift left] & {\Sigma'}^{\Phi}_{\pi} \arrow[l, shift left]
\end{tikzcd}\ .
\]

\msp{page 346}
% marge : trait manuscrit (reproduit en image) en regard de la Proposition 4.13.25
\textbf{Proposition 4.13.25.} \textit{Let $\mathcal{M}$, $\pi$, $\Phi$ be as in 4.13.23, with $I_{\alpha_\pi}\in\Phi$ ($\alpha_\pi$ being the smallest ordinal of cardinality $\Phi$} \texttt{[rather} $\pi$\texttt{]}\textit{). Then the set ${\Sigma'}^{\Phi}_{\pi}$ of subcategories of $\mathcal{M}$ $\Phi$-$\pi$-accessible in $\mathcal{M}$ is stable under intersection of families of cardinality $\leqslant\pi$} (\footnote{If $\pi$ is not a limit cardinal. Otherwise, one must say: families of cardinality $<\pi$. Cf. 4.13.10.3, where I give a slightly more general corrected statement.}).

(\textbf{N.B.} The bijection $\Sigma^{\Phi}_{\pi}\xrightarrow{\ \sim\ }{\Sigma'}^{\Phi}_{\pi}$ of 4.13.24 respects the natural order relations of inclusion, hence transforms intersections into intersections.)

Seen as in 4.13.7 (which is contained in the present statement), using lemma 4.13.10.

\textbf{4.13.26.} Let $\mathcal{M}$, $\pi$, $\Phi$ be as in 4.13.23, and let $\underline{S}$ be a \textit{strictly full} subcategory of $\mathcal{M}$, $f:\underline{S}\longrightarrow\mathcal{M}$ the inclusion functor. The condition that $\underline{S}$ be $\Phi$-$\pi$-accessible in $\mathcal{M}$ breaks up into:
\begin{itemize}[leftmargin=2.5em]
\item[1)] $\pi$ is useful for $\mathcal{M}$ \textit{and} for $\underline{S}$.
\item[2)] $\mathcal{M}$ and $\underline{S}$ are stable under $\varinjlim$ of type $\Phi$.
\item[3)] $f$ commutes with the said limits.
\item[4)] $f$ is $\pi$-accessible.
\item[5)] $f(\underline{S}_\pi)\subset\mathcal{M}_\pi$.
\end{itemize}

\printedpage{132}
% head: A. Grothendieck : Dérivateurs, Ch. XVIII, Édition des Universités de Montpellier II et Paris VII
% foot: Version 17 décembre 2025
\msp{page 347}
Let then $\pi'>\pi$ be another cardinal, when can we be sure that if $\underline{S}$ is $\Phi$-$\pi$-accessible in $\mathcal{M}$, it is also $\Phi$-$\pi'$-accessible\,? Conditions 2) 3) 4) are visibly stable under passage from $\pi$ to $\pi'$. Conditions 1) and 5) remain. We know that if $\pi'^{\pi}=\pi'$, then condition 1) remains valid for $\pi'$. More generally, if $\pi_0$ is a cardinal $\leqslant\pi$ such that $\pi_0$ is useful for $\mathcal{M}$ and for $\underline{S}$, then if $\pi'>\pi_0$ is such that $\pi'^{\pi_0}=\pi'$, then $\pi'$ is also useful for $\mathcal{M}$ and for $\underline{S}$. If we admit the GCH (generalized continuum hypothesis), this is so for every $\pi'>\pi_0$, hence for every $\pi'>\pi$.

Condition 5) remains, for which an answer is given by 4.12.1 (p. 273). If $\pi_0\leqslant\pi$ is chosen so that $f$ is $\pi_0$-accessible (e.g. $\pi_0=\pi$), and if
\[
\begin{aligned}
\pi'_1 &= \cardred\underline{S}_{\pi_0} \qquad\quad (\,\leqslant\cardred\underline{S}(\pi)\,)\\
\pi''_1 &= \sup(\pi'_1,\pi)\ ,
\end{aligned}
\]
we see that if $\pi'\geqslant\pi''_1$ (i.e. $\pi'\geqslant\pi'_1,\pi$) is such that $\pi'^{\pi_0}=\pi'$, then we do indeed have
\[ f(\underline{S}_{\pi'})\ \subset\ \mathcal{M}_{\pi'}\ . \]

\msp{page 348}
If then in addition $\pi'>\pi\geqslant\pi_0$ (as assumed at the start), and if we assume the GCH, then the condition $\pi'^{\pi_0}=\pi'$ is automatic, and only the condition $\pi'\geqslant\pi''_1$ remains, i.e. in the end $\pi'\geqslant\cardred\underline{S}_{\pi_0}$. To sum up:

\textbf{Proposition 4.13.27.} \textit{Under the conditions of 4.13.23 (without condition b) on $\Phi$), let $\underline{S}$ be a subcategory of $\mathcal{M}$ which is $\Phi$-$\pi$-adapted. Let $\pi_0\leqslant\pi$ be a cardinal which is useful for $\underline{S}$ and for $\mathcal{M}$, and such that $f:\underline{S}\longrightarrow\mathcal{M}$ is $\pi_0$-accessible, i.e. such that $\underline{S}$ is stable in $\mathcal{M}$ under $\varinjlim$ large relative to $\pi_0$ (one can e.g. take $\pi_0=\pi$). Then for every cardinal $\pi'>\pi$ such that}
\begin{itemize}[leftmargin=2.5em]
\item[\textit{a)}] \textit{$\pi'^{\pi_0}=\pi'$ and}
\item[\textit{b)}] \textit{$\pi'\geqslant\cardred\underline{S}_{\pi_0}$,}
\end{itemize}
\textit{$\underline{S}$ is $\Phi$-$\pi'$-accessible in $\mathcal{M}$. Thus if $\pi'>\pi$ is such that $\pi'^{\pi}=\pi'$, $\pi'>\cardred\underline{S}_\pi$, then $\underline{S}$ is $\Phi$-$\pi'$-accessible. (It suffices e.g. to take $\pi'=2^c$, with $c\geqslant\sup(\pi,\cardred\underline{S}(\pi))$.) If we admit the GCH, it therefore suffices to have $\pi'>\sup(\pi,\cardred\underline{S}(\pi))$.}

\msp{page 349}
\textbf{Proposition 4.13.28} (\footnote{To be moved further down, after 4.13.31, which is the first important application of 4.13.22, whereas 4.13.28 is proved without having to use 4.13.22 or its variants.}). \textit{Let $\mathcal{M}$, $\pi$, $\Phi$ be as in 4.13.23. Suppose in addition that $J\in\Phi$ implies that $J$ is} filtering. \textit{Then the set $\tilde{\Sigma}^{\Phi}_{\pi}$ of $\Phi$-$\pi$-accessible subsets of $\mathcal{M}$ is stable under finite unions. More generally, let $\lambda$ be a cardinal $<\pi$ such that $\forall\ J\in\Phi$, $J$ is large relative to $\lambda$. Then the set $\tilde{\Sigma}^{\Phi}_{\pi}$ is stable under unions $(S_r)_{r\in A}$, indexed by sets $A$ of cardinality $\leqslant\lambda$.}

\textbf{Corollary 4.13.29.} \textit{The set of accessible subsets of} $\Ob\mathcal{M}$ \textit{is stable under ``small'' unions -- i.e. if $(S_r)_{r\in A}$ is a small family of accessible subsets of} $\Ob\mathcal{M}$, \textit{then $\bigcup_{r\in A}S_r$ is also accessible.}

The proof of 4.13.28 rests on the following lemma (I do not spell out by the proof): % [sic] par la démonstration -> pas la démonstration, sense of the sentence

\printedpage{133}
% head: A. Grothendieck : Dérivateurs, Ch. XVIII, Édition des Universités de Montpellier II et Paris VII
% foot: Version 17 décembre 2025
\msp{page 350}
\textbf{Lemma 4.13.30.} \textit{Let $I$ be a category large relative to $\pi$, and $(J_r)_{r\in A}$ a family of full subcategories such that $\card A\leqslant\pi$, such that $\bigcup_{r\in A}\Ob J_r$ is cofinal in $\Ord(I)$, i.e. for every $i\in\Ob I$ $\exists\ j\in\bigcup_{r\in A}\Ob J_r$, and an arrow $i\longrightarrow j$. Then there exists $r\in A$ such that $J_r$ is cofinal in $I$.}

\textbf{Theorem 4.13.31.} \textit{Let $\mathcal{M}$, $\mathcal{N}$ be accessible categories, $f:\mathcal{M}\longrightarrow\mathcal{N}$ an accessible functor, $T\subset\Ob\mathcal{N}$ a subset of $\mathcal{N}$, $\underline{T}\subset\mathcal{N}$ the corresponding strictly full subcategory of $\mathcal{N}$, $S=f^{-1}(T)$, $\underline{S}=f^{-1}(\underline{T})$ the associated full subcategory of $\mathcal{M}$. If $T$ is accessible, then $S$ is accessible\,; in other words, if $\underline{T}$ is accessible in $\mathcal{N}$, $\underline{S}$ is accessible in $\mathcal{M}$.}

This follows from the more precise statement:

\textbf{Corollary 4.13.32.} \textit{Let $\pi_0$ and $\pi>\pi_0$ be cardinals having the following properties:}

\msp{page 351}
\begin{itemize}[leftmargin=2.5em]
\item[\textit{a)}] \textit{$\underline{T}$ is $\pi_0$-$\pi$-accessible in $\mathcal{N}$ (4.13.23).}
\item[\textit{b)}] \textit{$\pi_0$ is useful for $\mathcal{M}$ (it is for $\mathcal{N}$, because of a)).}
\item[\textit{c)}] \textit{$f(\mathcal{M}_\pi)\subset\mathcal{N}_\pi$.}
\item[\textit{d)}] \textit{$\mathcal{M}$ satisfies $L_{\pi_0}$, and $f$ is $\pi_0$-accessible.}
\end{itemize}
\textit{(The existence of such a pair $(\pi_0,\pi)$ follows at once from 4.13.16 and from 4.12.1.) Then $\underline{S}$ is $\pi_0$-$\pi$-accessible in $\mathcal{N}$.}

Since $f$ commutes with $\varinjlim$ large relative to $\pi_0$ by virtue of d), and $\underline{T}$ is stable under these limits by a) (cf. 4.13.23, 3°), the same is therefore true of $\underline{S}$, i.e. 3° of definition 4.13.23 is satisfied. On the other hand, condition 1° ($\pi$ useful for $\mathcal{M}$, $\mathcal{M}$ stable under $\varinjlim$ large relative to $\pi_0$ and of cardinality $\leqslant\pi$) is verified by b) and d). Thus it remains to verify condition 2°, i.e. that $\underline{S}$ is $\pi$-accessible. Let us put
\[
\begin{aligned}
\underline{T}(\pi) &= \underline{T}\cap\mathcal{N}_\pi\\
\underline{S}(\pi) &= \underline{S}\cap\mathcal{M}_\pi\ \ \underbrace{\ \subset\ }_{\text{thanks to c)}}\ \ f^{-1}(\underline{T}(\pi))\ .
\end{aligned}
\]
Since $\underline{S}$ is stable in $\mathcal{M}$ under $\varinjlim$ large relative to $\pi_0$ (which exist in $\mathcal{M}$ by d), and a fortiori under $\varinjlim$ large relative to $\pi$, we have a fully faithful canonical functor
\[ \Ind_\pi\underline{S}(\pi)\ \longrightarrow\ \underline{S}\ . \]

\msp{page 352}
Besides, $\underline{T}$ being strictly full in $\mathcal{N}$, $\underline{S}$ is strictly full in $\mathcal{M}$, and it remains to verify that the preceding functor is essentially surjective. Let then $x\in\Ob\underline{S}$. We will have, thanks to b),
\[ x\ =\ \varinjlim_{I}x_i\ ,\qquad\quad x_i\in\mathcal{M}_\pi,\ I=\mathcal{M}_\pi/x\ . \]
Besides, the preceding direct system is $\pi_0$-$\pi$-adapted, i.e. (p. 339):
\begin{itemize}[leftmargin=2.5em]
\item[1)] $I$ equivalent to a small category, $I$ large relative to $\pi$ (follows from the fact that $\pi$ is useful for $\mathcal{M}$, cf.\,\cercle{b}\,).
\end{itemize}

\end{document}
